% 核心结论：理想变压器的开路、短路与直流边界
% 图原型：纯几何／对比。物理网络按本题节点、端口和符号构建。
\documentclass[tikz,border=7pt]{standalone}
\usepackage{fontspec,xeCJK,amsmath,amssymb}
\IfFontExistsTF{Times New Roman}{\setmainfont{Times New Roman}}{\setmainfont{TeX Gyre Termes}}
\IfFontExistsTF{Heiti SC}{\setCJKmainfont[AutoFakeBold=1.5]{Heiti SC}}{\setCJKmainfont[AutoFakeBold=1.5]{FandolHei-Regular}}
\usetikzlibrary{arrows.meta,calc,angles,quotes}
\definecolor{ink}{HTML}{243745}
\definecolor{blue}{HTML}{275DAB}
\definecolor{green}{HTML}{177C63}
\definecolor{red}{HTML}{BD4935}
\definecolor{purple}{HTML}{715393}
\begin{document}
\begin{tikzpicture}[font=\fontsize{9.5}{12}\selectfont,text=ink,
 v/.style={line width=1.1pt,-{Stealth[length=2.4mm,width=1.5mm]}},
 axis/.style={ink!65,line width=.45pt,-{Stealth[length=1.7mm,width=1.1mm]}},
 aux/.style={ink!40,densely dashed,line width=.5pt},
 lab/.style={fill=white,inner sep=2pt},
 card/.style={draw=blue!25,fill=blue!3,rounded corners=3pt,text width=6cm,align=left,inner sep=10pt}]
\path[use as bounding box] (-8,-3.35) rectangle (8,3.65);
\node[font=\fontsize{10.5}{13}\selectfont\bfseries,text width=15cm,align=center] at (0.00000,3.22000) {M6-H28　理想变压器的开路、短路与直流边界};
\begin{scope}[shift={(-4,0)}]
\draw[ink,thick] (-2.60000,1.50000) -- (-2.60000,0.40000);
\draw[ink,thick] (-2.60000,-0.40000) -- (-2.60000,-1.50000);
\draw[ink,thick] (-2.6,0) circle (.4);
\node[] at (-2.60000,0.00000) {$\sim$};
\node[] at (-3.22000,0.00000) {$U_1$};
\draw[ink,thick] (-2.60000,1.50000) -- (-0.90000,1.50000);
\draw[ink,thick] (-2.60000,-1.50000) -- (-0.90000,-1.50000);
\draw[ink,thick] (-0.90000,1.50000) -- (-0.90000,0.42000);
\draw[ink,thick] (-0.90000,-0.42000) -- (-0.90000,-1.50000);
\begin{scope}[shift={(-0.90000,0.00000)},rotate=-90.0]
\draw[ink,thick] plot[domain=-0.42:0.42,samples=70] (\x,{.20*abs(sin(720*(\x+0.42)/0.84))});
\end{scope}
\node[lab,anchor=east] at (-1.32000,0.00000) {$N_1$};
\draw[ink,thick] (0.90000,1.50000) -- (0.90000,0.42000);
\draw[ink,thick] (0.90000,-0.42000) -- (0.90000,-1.50000);
\begin{scope}[shift={(0.90000,0.00000)},rotate=-90.0]
\draw[ink,thick] plot[domain=-0.42:0.42,samples=70] (\x,{.20*abs(sin(720*(\x+0.42)/0.84))});
\end{scope}
\node[lab,anchor=west] at (1.32000,-0.65000) {$N_2$};
\draw[ink!50,line width=1.4pt] (0.00000,-1.00000) -- (0.00000,1.00000);
\draw[ink!50,line width=1.4pt] (0.20000,-1.00000) -- (0.20000,1.00000);
\draw[ink,thick] (0.90000,1.50000) -- (2.60000,1.50000);
\draw[ink,thick] (2.60000,-1.50000) -- (0.90000,-1.50000);
\draw[ink,thick] (2.60000,1.50000) -- (2.60000,0.42000);
\draw[ink,thick] (2.60000,-0.42000) -- (2.60000,-1.50000);
\begin{scope}[shift={(2.60000,0.00000)},rotate=-90.0]
\draw[ink,thick,fill=white] (-0.42,-.14) rectangle (0.42,.14);
\end{scope}
\node[lab,anchor=west] at (3.02000,0.00000) {$Z_L$};
\fill[ink] (-0.90000,0.80000) circle (1.6pt);
\fill[ink] (0.90000,0.80000) circle (1.6pt);
\draw[v,blue] (-2.35000,1.50000) -- (-1.50000,1.50000);
\node[lab,text=blue,anchor=center] at (-1.92500,1.82000) {$I_1$};
\draw[v,green] (1.20000,1.50000) -- (2.00000,1.50000);
\node[lab,text=green,anchor=center] at (1.60000,1.82000) {$I_{out}$};
\node[] at (0.00000,-2.40000) {上端为同名端；次级负载电流从同名端流出。};
\end{scope}
\begin{scope}[shift={(4,0)}]
\node[card] at (0.00000,1.15000) {$Z_L\to\infty:\ I_1\to0$};
\end{scope}
\begin{scope}[shift={(4,0)}]
\node[card] at (0.00000,-1.15000) {$Z_L=0\text{ 与理想恒压驱动冲突}$};
\end{scope}
\node[text width=15cm,align=center] at (0.00000,-3.00000) {同名端、电流方向与折算侧统一；有限励磁和损耗不能混入理想模型。};

\end{tikzpicture}\end{document}